This section sets out the gaps left open by prior work and summarises how this
study addresses them.

\subsection{Research Gap}

Despite substantial progress in data-driven SOH prediction, three critical gaps remain unaddressed in the existing literature:

\begin{enumerate}
  \item \textbf{Cell-to-cell variability is not modelled.}
        Published benchmarks typically use hand-picked cells with moderate,
        smooth degradation trajectories.
        Cell-to-cell variance, to which manufacturing tolerance is a likely
        contributor, can produce 10$\times$
        differences in degradation rate under identical conditions, is
        systematically understudied.
        Most transfer and fine tuning methods adapt to a target cell either by
        aligning feature distributions across batteries at the population
        level~\cite{ye2021state} or by fine tuning a shared model on the target
        cell's own recent data~\cite{ma2022real}.
        Neither extracts an explicit early-life shape signature that locates a
        cell within the manufacturing distribution of the pretrained population
        before adaptation, so a cell whose early trajectory is ambiguous
        (Section~\ref{sec:variability}) cannot be distinguished from an average
        cell.

  \item \textbf{The standard benchmark protocol has never been validated.}
        Nearly all published comparisons, including our own prior
        work~\cite{ponnambalam2024battery,ponnambalam2026delstm}, rank methods
        by one-step-ahead RMSE on small cylindrical 18650 cells (with NASA cells
        B05, B06, B07, and B18 the de facto standard since
        2007~\cite{saha2007battery,xu2025review}).
        Whether this protocol can actually distinguish methods, that is,
        whether the reported differences exceed trivial-baseline performance
        and the run-to-run variance of training itself, has never been
        tested, and neither has its behaviour on large-format
        ($\geq$10\,Ah) cells under industrially representative conditions.

  \item \textbf{Uncertainty calibration is chemistry-dependent.}
        Bayesian approximations and conformal prediction, including MC-Dropout
        and Bayes by Backprop (BBB), provide well-studied uncertainty
        estimates~\cite{gal2016dropout,angelopoulos2021gentle} and have been
        applied to battery SOH~\cite{ke2024bayesian}, but typically within a
        single chemistry.
        Whether these estimates remain calibrated under a chemistry mismatch
        between pretraining and deployment data has not been investigated, a
        critical gap for practical BMS deployment where pretraining data may be
        available only for a different battery chemistry.
\end{enumerate}

This study addresses these three gaps using the SIT
dataset~\cite{ponnambalam2026sitdata}, our previously published large-format
campaign of 20 LiFePO$_4$ prismatic cells (50\,Ah, 3.2\,V) cycled over 18 months
under three operating conditions. As the baseline we use our previously published
DE-LSTM~\cite{ponnambalam2026delstm}, a differential-evolution-optimised LSTM with
leading accuracy on the NASA 18650 benchmark, alongside five further deep
baselines and two trivial baselines. Section~\ref{sec:problem} quantifies how
severely the first two gaps, cell-to-cell variability and benchmark saturation,
manifest on these cells; Sections~\ref{sec:method} and~\ref{sec:results} present
and validate the proposed hybrid, whose value lies in deployment capabilities
rather than one-step rank; and a cross-dataset evaluation establishes the
chemistry-alignment effect, in which the model's uncertainty estimates stay
calibrated only when pretraining and deployment chemistries match.

\subsection{Contributions}
This paper makes the following key contributions:
\begin{enumerate}
  \item \textbf{Large-format benchmark with trivial-baseline control}: the
        first systematic evaluation of six deep single-cell predictors,
        including our prior DE-LSTM, on 50\,Ah LiFePO$_4$ prismatic cells,
        together with linear-autoregressive and curve-fitting controls that
        expose the saturation of the one-step protocol
        (Section~\ref{sec:baseline}).

  \item \textbf{Cell-to-cell variability quantification}: Across ten nominally
        identical G1 cells, we show that although each cell's own degradation
        trajectory is highly predictable, the cells differ so much from one
        another that this spread between cells, not the algorithm, sets the
        error floor.
        A variance decomposition across six predictors attributes 25\% to 58\%
        of error variance to the cell and only ${\sim}$15\% to the method,
        and cross-cell pooling exhibits a two-sided failure mode: it rescues
        the fast-degrading cells but worsens others by up to 2$\times$
        (Section~\ref{sec:variability}).

  \item \textbf{Benchmark-saturation finding}: We show that the field's
        standard one-step-ahead protocol cannot rank methods on data-rich
        large-format cells: a linear autoregressor attains the lowest error of
        all evaluated models, inter-method differences fall within
        training-run variance, and apparent cross-condition effects are
        dominated by the quality of individual training runs
        (Sections~\ref{sec:baseline} and~\ref{sec:brittleness}).

  \item \textbf{Hybrid transfer model with ablation}: We propose a three-phase
        model (MIT LFP pretraining, convolutional fingerprint encoder,
        MC-Dropout uncertainty) and conduct a full ablation study isolating
        each component's contribution.
        The hybrid predicts new cells from 30 cycles with no target-cell
        fine-tuning (23\% better than scratch) and achieves a 51\%
        mean RMSE improvement over scratch with fine tuning, on all 17 of 17
        test cells; it is the only evaluated method that quantifies
        uncertainty, with prediction intervals that are recalibrated by
        grouped conformal inference (Section~\ref{sec:solution}).
        We additionally characterise pretraining-run variance, show that
        source-domain criteria cannot detect poorly transferring pretraining
        runs, and introduce a calibration-cell checkpoint-selection protocol
        that makes the pipeline reliably reproducible
        (Section~\ref{sec:selection}).

  \item \textbf{Chemistry-alignment effect with multi-dataset validation}:
        We demonstrate that MC-Dropout uncertainty calibration transfers between
        datasets only when pretraining and deployment chemistries are matched
        (LFP$\rightarrow$LFP: 80.8\% coverage; LFP$\rightarrow$LCO: 43.6\%
        coverage).
        A ten-scenario cross-dataset experiment spanning three LFP datasets
        (MIT, SIT, LISHEN 40\,Ah) and one LCO dataset confirms this effect
        and shows that fine tuning is the dependable adaptation mechanism
        across chemistries, while fingerprint conditioning without fine-tuning
        provides its distinctive value before any adaptation is possible, yielding a
        practical guideline for transfer learning in BMS applications
        (Sections~\ref{sec:cross_dataset} and~\ref{sec:discussion}).

  \item \textbf{Early screen for manufacturing-defective cells}: We show that
        no scalar health indicator separates the defective cells within a group,
        yet the early-cycle discharge voltage-curve descriptor $\Delta Q(V)$
        detects eventual early death at an AUC of 0.78 on the SIT cells,
        reproduced on the MIT corpus, about 50 cycles before the capacity
        trajectory reveals the defect (Section~\ref{sec:screening}).

  \item \textbf{Full-coverage probabilistic RUL forecasting}: We extend the
        hybrid from one-step SOH to remaining-useful-life forecasting with a
        direct multi-horizon decoder and a 155-cell multi-dataset LFP foundation
        corpus, forecasting every cell that reaches end of life, including the
        catastrophic cells that recursive baselines abstain on, at a median
        error of 49 cycles at the midlife anchor with conformally calibrated
        intervals (Section~\ref{sec:rul}).
\end{enumerate}
